%%%%%%%%%%%%%%%%%%%%%%%%%%%%%%%%%%%%%%%%%%%%%%%%%%%%%%%%%%%%
\section{Methods}
\label{sec:methods}
%%%%%%%%%%%%%%%%%%%%%%%%%%%%%%%%%%%%%%%%%%%%%%%%%%%%%%%%%%%%

Both \molm~\citep{ross2022largescalechemicallanguagerepresentations} and \smi~\citep{soares2024smi} use 12-layer, 12-head transformer encoders with hidden dimension $d=768$. We run the same two experiments on both models to compare whether their learned representations encode molecular structure in similar ways. The attention-distance experiment tests whether attention heads directly track 3D spatial proximity, while the probing experiment tests whether atom-level chemical properties are linearly recoverable from hidden states. Throughout, we use RDKit~\cite{rdkit}, an open-source cheminformatics toolkit, to access molecular properties and generate 3D conformers.

\subsection{Attention-Distance Correlation.} For each attention head in every layer, we compare attention weights between \textbf{heavy} (non-Hydrogen) atom pairs $A_{ij}^{(l,h)}$ to their corresponding 3D Euclidean distances $D_{ij}$ (measured in \AA).\footnote{Since SMILES strings do not contain explicit 3D geometric information, we use ETKDGv3 and MMFF94 (modules from RDKit) to generate a 3D coordinate $\mathbf{r}_i=(x_i,y_i,z_i)$ for each atom. This procedure is used only as a preprocessing step; no additional model training is performed.}.

For each layer-head pair, we compute the Spearman correlation $\rho$ between attention weights and inverse distance separately for every molecule, then average the resulting correlations across the sample. More detailed setup is in Appendix~\ref{app:attention-appendix}.

We additionally stratify atom pairs by distance into \emph{short} ($D_{ij}\le2$\,\AA), \emph{medium} ($2<D_{ij}\le4$\,\AA), and \emph{long} ($D_{ij}>4$\,\AA) interactions, and repeat the analysis within each group.

\subsection{Probing Chemical Properties}
\para{Probing.}
Let $\mathbf{h}^{(i)}_l \in \mathbb{R}^d$ denote the residual-stream hidden state at the $i$-th atom-token position at layer $l$, where $d{=}768$ and $l \in \{0, 1, \dots, 12\}$ (layer $0$ is the embedding output). We fit a linear probe $\mathbf{h} \mapsto \mathbf{W}\mathbf{h} + \mathbf{b}$ with $\mathbf{W} \in \mathbb{R}^{C_\tau \times d}$ and $\mathbf{b} \in \mathbb{R}^{C_\tau}$, where $C_\tau$ is the number of classes for task $\tau$. We compare the linear probe to two references: (i) a \emph{random-input} linear probe on a Gaussian feature matrix $X_{ij} \sim \mathcal{N}(0, 1)$ in place of the real hidden states (controls for head capacity) and (ii) a \emph{majority-vote} baseline that always predicts the most frequent training-set class. We report test accuracy and macro-F1.

\para{Causal Intervention.} For a given task, choice of source class $c$, target class $c'$, and steer layer $\ell$, we shift every atom-token position whose baseline prediction is $c$ by $\alpha \cdot (\mathbf{W}_{c'}-\mathbf{W}_{c})/\boldsymbol{\sigma}$ at layer $\ell$, then run the rest of the encoder and collect the layer-$12$ prediction. \emph{Intervention success rate} is the fraction of those positions whose post-intervention layer-$12$ prediction equals $c'$. We run a separate single-layer intervention at each candidate $\ell \in \{1, \ldots, 12\}$, using the probe trained at that layer's hidden states to define both the source-class detector and the steering direction, and report the success rate at the best-performing $\alpha$ on held-out molecules. More detailed setup is in Appendix~\ref{app:steering}.

\para{Tasks.} We probe eight per-atom RDKit-derived properties on QM9~\citep{ramakrishnan2014quantum}: atom type ($C_\tau{=}4$), hybridization ($C_\tau{=}3$), chiral tag ($C_\tau{=}3$), aromaticity ($C_\tau{=}2$), ring membership ($C_\tau{=}2$), degree ($C_\tau{=}4$), formal charge ($C_\tau{=}3$), and total valence ($C_\tau{=}4$). Integer-valued properties (degree, valence) are treated as discrete classes rather than ordinal regression. Atom-token alignment, train/test split, optimisation hyperparameters, and standardisation are detailed in Appendix~\ref{app:training}.



% \subsection{Experiment 2: Linear Probing}

% We test where chemically meaningful atom-level information emerges across model
% layers by training simple linear probes on hidden states. For each SMILES string \textbf{(1000 random selected from QM9 dataset)},
% we construct an RDKit molecule and extract atom-level labels including atom type,
% hybridization, aromaticity, ring membership, chirality, degree, formal charge, and
% total valence. We align SMILES tokens to atoms and extract the corresponding hidden
% state $h_i^{(\ell)} \in \mathbb{R}^d$ for each atom $i$ at each layer $\ell$.

% For each property, we learn a linear mapping from hidden states to labels. Since these are categorical tasks, we use multinomial logistic regression:
% \[
% \hat{y}_i = \mathrm{softmax}(W^{(\ell)} h_i^{(\ell)} + b^{(\ell)}),
% \]
% where $W^{(\ell)}$ and $b^{(\ell)}$ are layer-specific parameters. 

% We split the dataset at the molecule level (80/20) so that all atoms from a given molecule
% appear only in either the train or test set. Probe performance is evaluated using
% accuracy and weighted F1, and compared against a most-frequent-class dummy baseline.

% The main difference between models is how hidden states are extracted. For
% MoLFormer-XL, hidden states $\{h^{(\ell)}\}_{\ell=0}^L$ are returned directly via
% the HuggingFace API using \texttt{output\_hidden\_states=True}. For
% SMI-TED$_{289\mathrm{M}}$, hidden states are captured by registering forward hooks
% on the token embedding layer and each encoder block. After extraction, the probing
% pipeline is identical across models: atom-token alignment, molecule-level
% train/test splitting, and layer-wise linear probing.

% [Need to do for molecular level properties as well.]

% [running into issue with inherent biases in QM9 dataset.]

% \subsection{Experiment 3: Representation Geometry via PCA and UMAP}

% We investigate the structure of learned atom-level representations by projecting
% hidden states into low-dimensional spaces using principal component analysis
% (PCA) and uniform manifold approximation and projection (UMAP). This experiment
% provides a qualitative and geometric complement to the linear probing analysis,
% allowing us to visualize how chemical properties are organized in representation
% space across layers.

% \paragraph{Atom-level representation extraction.}
% For each SMILES string, we construct the corresponding RDKit molecule and align
% SMILES tokens to heavy atoms using a tokenizer-based mapping. For each atom $i$
% and layer $\ell$, we extract the hidden state
% $h_i^{(\ell)} \in \mathbb{R}^d$ from the encoder. We analyze representations at
% multiple layers, including the embedding layer ($\ell=0$), intermediate layers,
% and the final layer.



% \paragraph{Dimensionality reduction.}
% For each layer $\ell$, we standardize hidden states across all atoms and apply
% PCA to obtain a two-dimensional projection:
% \[
% z_i^{(\ell)} = \mathrm{PCA}_2\bigl(h_i^{(\ell)}\bigr).
% \]
% We additionally apply UMAP with cosine distance to capture potential
% nonlinear structure in the representation space.

% \paragraph{Random baseline.}
% To assess whether observed structure arises from training rather than architectural
% bias, we repeat the entire analysis using a randomly initialized encoder with
% identical architecture. This provides a baseline for comparison against the
% pretrained model.

% \paragraph{Objective.}

% [still in the works]

% This experiment is designed to answer the following questions:
% \begin{itemize}
%     \item Do atom-level representations cluster according to chemical properties?
%     \item At which layers does clustering emerge?
%     \item Does the learned representation exhibit nonlinear organization not
%     captured by linear probes?
%     \item How does pretrained structure compare to random initialization?
% \end{itemize}

% [need to analyze other lower variance directions because information about properties could be stored there: LDA]

% \subsection{Experiment 4: Supervised Representation Structure via LDA}

% While PCA provides an unsupervised view of representation geometry, it may fail
% to reveal chemically meaningful structure if such information lies in
% low-variance directions. To address this limitation, we perform a supervised
% dimensionality reduction analysis using linear discriminant analysis (LDA),
% which explicitly identifies directions that maximize class separability.

% \paragraph{Atom-level representation extraction.}
% As in previous experiments, we extract atom-level hidden states
% $h_i^{(\ell)} \in \mathbb{R}^d$ for each atom $i$ and layer $\ell$, using
% token-to-atom alignment based on SMILES tokenization and RDKit molecule
% construction. We analyze representations across multiple layers, including the
% embedding layer and intermediate and final transformer layers.

% \paragraph{LDA projection.}
% For each property and each layer $\ell$, we standardize hidden states and fit a
% linear discriminant analysis model:
% \[
% z_i^{(\ell)} = \mathrm{LDA}_k\bigl(h_i^{(\ell)}\bigr),
% \]
% where $k = \min(2, C - 1)$ and $C$ is the number of classes. For binary
% properties, this yields a one-dimensional projection, which we augment with
% random jitter for visualization. For multi-class properties, we use the top two
% discriminant components.

% \paragraph{Class balancing.}
% To mitigate class imbalance---particularly the dominance of carbon atoms---we
% construct a balanced subset of atoms by subsampling frequent classes. This
% ensures that LDA projections are not dominated by a single majority class and
% that separability reflects meaningful chemical distinctions.

% \paragraph{Probe-direction comparison.}
% To further connect supervised projections with probing results, we compare LDA
% directions to linear probe directions obtained from logistic regression. For
% binary tasks, we project representations onto the probe weight vector; for
% multi-class tasks, we project onto the top singular vectors of the probe weight
% matrix. This provides a direct link between representation geometry and
% predictive features.

% \paragraph{Objective.}
% This experiment is designed to answer the following questions:
% \begin{itemize}
%     \item Do atom representations become linearly separable by chemical properties?
%     \item At which layers does class separation emerge?
%     \item How does supervised structure compare to unsupervised PCA?
%     \item Are discriminative directions aligned with linear probe features?
% \end{itemize}

% \subsection{Experiment 5: Attention Head Ablation.}
% To be described.
